\chapter{2008年湖南卷（理）}

\section{单选题}

\begin{problem}
复数 \(\left(\i - \frac{1}{\i}\right)^3\) 等于（\quad）

\choices
    {\(8\)}
    {\(-8\)}
    {\(8\i\)}
    {\(-8\i\)}
\end{problem}

\begin{answer}
D
\end{answer}

\begin{solution}
因为 \(\i^{-1}=-\i\)，所以
\[
\left(\i-\frac{1}{\i}\right)^3=(2\i)^3=8\i^3=-8\i
\]
故选 D.
\end{solution}

\begin{problem}
“\(\lvert x-1\rvert<2\) 成立”是“\(x(x-3)<0\) 成立”的（\quad）

\choices
    {充分不必要条件}
    {必要不充分条件}
    {充分必要条件}
    {既不充分也不必要条件}
\end{problem}

\begin{answer}
B
\end{answer}

\begin{solution}
设命题 \(\mathcal A\) 为 \(\lvert x-1\rvert<2\)，命题 \(\mathcal B\) 为 \(x(x-3)<0\). 由
\(\lvert x-1\rvert<2\Longleftrightarrow -1<x<3\)，\(x(x-3)<0\Longleftrightarrow 0<x<3\).
可知 \(\mathcal B\) 成立时 \(\mathcal A\) 必成立，但取 \(x=-\frac12\) 时 \(\mathcal A\) 成立而 \(\mathcal B\) 不成立. 因此，\(\mathcal A\) 是 \(\mathcal B\) 的必要不充分条件，故选 B.
\end{solution}

\begin{problem}
已知变量 \(x\)，\(y\) 满足条件 \(\begin{cases}
x \geqslant 1,\\
x - y \leqslant 0,\\
x + 2y - 9 \leqslant 0,
\end{cases}\) 则 \(x + y\) 的最大值是（\quad）

\choices
    {\(2\)}
    {\(5\)}
    {\(6\)}
    {\(8\)}
\end{problem}

\begin{answer}
C
\end{answer}

\begin{solution}
由约束条件得 \(x\geqslant1\)，\(y\geqslant x\)，且 \(y\leqslant\frac{9-x}{2}\). 可行域由直线 \(x=1\)、\(y=x\) 和 \(x+2y=9\) 围成，三个顶点分别为
\((1,1)\)，\((1,4)\)，\((3,3)\).
在三个顶点处，\(x+y\) 的值分别为 \(2\)、\(5\)、\(6\). 线性目标函数在多边形可行域上的最大值取在顶点处，所以最大值为 \(6\)，故选 C.
\end{solution}

\begin{problem}
设随机变量 \(\xi\) 服从正态分布 \(N(2,9)\)，若 \(P(\xi > c + 1) = P(\xi < c - 1)\)，则 \(c =\)（\quad）

\choices
    {\(1\)}
    {\(2\)}
    {\(3\)}
    {\(4\)}
\end{problem}

\begin{answer}
B
\end{answer}

\begin{solution}
正态分布 \(N(2,9)\) 的密度关于直线 \(x=2\) 对称. 因而
\[
P(\xi>c+1)=P\bigl(\xi<4-(c+1)\bigr)=P(\xi<3-c)
\]
题设还给出该概率等于 \(P(\xi<c-1)\). 正态分布的分布函数严格递增，所以
\[
3-c=c-1
\]
解得 \(c=2\)，故选 B.
\end{solution}

\begin{problem}
设有直线 \(m\)，\(n\) 和平面 \(\alpha\)，\(\beta\). 下列四个命题中，正确的是（\quad）

\choices
    {若 \(m \parallel \alpha, n \parallel \alpha\)，则 \(m \parallel n\)}
    {若 \(m \subset \alpha\)，\(n \subset \alpha\)，\(m \parallel \beta\)，\(n \parallel \beta\)，则 \(\alpha \parallel \beta\)}
    {若 \(\alpha \perp \beta\)，\(m \subset \alpha\)，则 \(m \perp \beta\)}
    {若 \(\alpha \perp \beta\)，\(m \perp \beta\)，\(m \not\subset \alpha\)，则 \(m \parallel \alpha\)}
\end{problem}

\begin{answer}
D
\end{answer}

\begin{solution}
命题 A 不成立. 例如取 \(\alpha:z=0\)，\(m:\{(t,0,1):t\in\R\}\)，\(n:\{(0,t,1):t\in\R\}\)，两直线都平行于 \(\alpha\)，但相交于 \((0,0,1)\)，所以不平行.

命题 B 不成立. 取 \(\alpha:z=0\)，\(\beta:x=0\)，并取 \(m:\{(1,t,0):t\in\R\}\)，\(n:\{(2,t,0):t\in\R\}\). 两直线都在 \(\alpha\) 内且都平行于 \(\beta\)，但 \(\alpha\) 与 \(\beta\) 相交.

命题 C 不成立. 当 \(\alpha\perp\beta\) 时，取 \(m=\alpha\cap\beta\)，则 \(m\subset\alpha\)，但 \(m\) 在 \(\beta\) 内，不能与 \(\beta\) 垂直.

对于命题 D，\(m\perp\beta\) 表明直线 \(m\) 的方向平行于平面 \(\beta\) 的法线. 由于 \(\alpha\perp\beta\)，平面 \(\beta\) 的法线平行于平面 \(\alpha\)，故 \(m\parallel\alpha\) 或 \(m\subset\alpha\). 题设又有 \(m\not\subset\alpha\)，所以 \(m\parallel\alpha\). 故选 D.
\end{solution}

\begin{problem}
函数 \( f(x) = \sin^2 x + \sqrt{3}\sin x\cos x \) 在区间 \(\left[\frac{\pi}{4},\frac{\pi}{2}\right]\) 上的最大值是（\quad）

\choices
    {\(1\)}
    {\(\frac{1 + \sqrt{3}}{2}\)}
    {\(\frac{3}{2}\)}
    {\(1 + \sqrt{3}\)}
\end{problem}

\begin{answer}
C
\end{answer}

\begin{solution}
利用二倍角公式，得
\[
f(x)=\frac{1-\cos 2x}{2}+\frac{\sqrt{3}}{2}\sin 2x
=\frac12+\sin\left(2x-\frac{\pi}{6}\right)
\]
当 \(x\in\left[\frac{\pi}{4},\frac{\pi}{2}\right]\) 时，\(2x-\frac{\pi}{6}\in\left[\frac{\pi}{3},\frac{5\pi}{6}\right]\)，其中包含正弦函数取得最大值 \(1\) 的点 \(\frac{\pi}{2}\). 因此 \(f(x)\) 的最大值为 \(\frac32\)，故选 C.
\end{solution}

\begin{problem}
设 \(D\)，\(E\)，\(F\) 分别是 \(\triangle ABC\) 的三边 \(BC\)，\(CA\)，\(AB\) 上的点，且 \(\overrightarrow{DC} = 2\overrightarrow{BD}\)，\(\overrightarrow{CE} = 2\overrightarrow{EA}\)，\(\overrightarrow{AF} = 2\overrightarrow{FB}\)，则 \(\overrightarrow{AD} + \overrightarrow{BE} + \overrightarrow{CF}\) 与 \(\overrightarrow{BC}\)（\quad）

\choices
    {反向平行}
    {同向平行}
    {互相垂直}
    {既不平行也不垂直}
\end{problem}

\begin{answer}
A
\end{answer}

\begin{solution}
由 \(\overrightarrow{DC}=2\overrightarrow{BD}\)，得 \(\overrightarrow{BD}=\frac13\overrightarrow{BC}\)，从而
\[
\overrightarrow{AD}=\overrightarrow{AB}+\frac13\overrightarrow{BC}
=\frac23\overrightarrow{AB}+\frac13\overrightarrow{AC}
\]
由 \(\overrightarrow{CE}=2\overrightarrow{EA}\) 及 \(\overrightarrow{CA}=\overrightarrow{CE}+\overrightarrow{EA}\)，得
\[
\overrightarrow{CE}=\frac23\overrightarrow{CA}
\]
因此
\[
\overrightarrow{BE}=\overrightarrow{BC}+\overrightarrow{CE}
=\left(\overrightarrow{AC}-\overrightarrow{AB}\right)-\frac23\overrightarrow{AC}
=\frac13\overrightarrow{AC}-\overrightarrow{AB}
\]
同样，由 \(\overrightarrow{AF}=2\overrightarrow{FB}\) 及 \(\overrightarrow{AB}=\overrightarrow{AF}+\overrightarrow{FB}\)，得
\[
\overrightarrow{AF}=\frac23\overrightarrow{AB}
\]
因此
\[
\overrightarrow{CF}=\overrightarrow{CA}+\overrightarrow{AF}
=-\overrightarrow{AC}+\frac23\overrightarrow{AB}
\]
所以
\[
\overrightarrow{AD}+\overrightarrow{BE}+\overrightarrow{CF}
=\frac13\overrightarrow{AB}-\frac13\overrightarrow{AC}
=-\frac13\overrightarrow{BC}
\]
该向量与 \(\overrightarrow{BC}\) 反向平行，故选 A.
\end{solution}

\begin{problem}
若双曲线 \(\frac{x^2}{a^2} - \frac{y^2}{b^2} = 1\)（\(a > 0\)，\(b > 0\)）上横坐标为 \(\frac{3a}{2}\) 的点到右焦点的距离大于它到左准线的距离，则双曲线离心率的取值范围是（\quad）

\choices
    {\((1,2)\)}
    {\((2, +\infty)\)}
    {\((1,5)\)}
    {\((5, +\infty)\)}
\end{problem}

\begin{answer}
B
\end{answer}

\begin{solution}
设双曲线的半焦距为 \(c\)，离心率为 \(e=\frac ca>1\)，则 \(b^2=a^2(e^2-1)\). 设题设点为 \(Q\)，其横坐标为 \(\frac{3a}{2}\)，由双曲线方程得
\[
y_Q^2=b^2\left(\frac{9}{4}-1\right)=\frac54b^2
\]
因此到右焦点 \((c,0)\) 的距离满足
\[
QF^2=\left(\frac{3a}{2}-c\right)^2+\frac54b^2
=\frac{a^2}{4}(3e-2)^2
\]
由于 \(e>1\)，所以 \(QF=\frac a2(3e-2)\). 左准线方程为 \(x=-\frac{a^2}{c}\)，故点 \(Q\) 到左准线的距离为
\[
Q\ell=\frac{3a}{2}+\frac{a^2}{c}=\frac a{2e}(3e+2)
\]
由题设不等式
\[
\frac a2(3e-2)>\frac a{2e}(3e+2)
\Longleftrightarrow 3e^2-5e-2>0
\Longleftrightarrow (3e+1)(e-2)>0
\]
结合 \(e>1\)，得 \(e>2\)，故选 B.
\end{solution}

\begin{problem}
长方体 \(ABCD - A_{1}B_{1}C_{1}D_{1}\) 的 \(8\) 个顶点在同一球面上，且 \(AB = 2\)，\(AD = \sqrt{3}\)，\(AA_{1} = 1\)，则顶点 \(A\)，\(B\) 间的球面距离是（\quad）

\choices
    {\(2\sqrt{2}\pi\)}
    {\(\sqrt{2}\pi\)}
    {\(\frac{\sqrt{2}\pi}{2}\)}
    {\(\frac{\sqrt{2}\pi}{4}\)}
\end{problem}

\begin{answer}
C
\end{answer}

\begin{solution}
长方体的体对角线长为
\[
\sqrt{AB^2+AD^2+AA_1^2}=\sqrt{4+3+1}=2\sqrt2
\]
该体对角线是外接球的直径，所以球半径为 \(R=\sqrt2\). 设球心为 \(O\)，\(\theta=\angle AOB\) 为小于等于 \(\pi\) 的圆心角，则弦长公式给出
\[
2=AB=2R\sin\frac\theta2=2\sqrt2\sin\frac\theta2
\]
于是 \(\theta=\frac\pi2\). 球面上两点间的最短距离为对应大圆弧长，因此
\[
R\theta=\sqrt2\cdot\frac\pi2=\frac{\sqrt2\pi}{2}
\]
故选 C.
\end{solution}

\begin{problem}
设 \( [x] \) 表示不超过 \(x\) 的最大整数（如 \([2]=2\)，\(\left[\frac{5}{4}\right]=1\) ）. 对于给定的 \( n \in \mathbf{N}^{*} \)，定义 \( C_{x}^{n} = \frac{n(n-1)\cdots(n-[x]+1)}{x(x-1)\cdots(x-[x]+1)} \)， \( x \in [1, +\infty) \)，则当 \( x \in \left[\frac{3}{2}, 3\right) \) 时，函数 \( C_{x}^{8} \) 的值域是（\quad）

\choices
    {\(\left[\frac{16}{3}, 28\right]\)}
    {\(\left[\frac{16}{3}, 56\right)\)}
    {\(\left(4, \frac{28}{3}\right) \cup [28, 56)\)}
    {\(\left(4, \frac{16}{3}\right] \cup \left(\frac{28}{3}, 28\right]\)}
\end{problem}

\begin{answer}
D
\end{answer}

\begin{solution}
当 \(\frac32\leqslant x<2\) 时，\([x]=1\)，所以
\(C_x^8=\frac8x\)，\(4<\frac8x\leqslant\frac{16}{3}\).
当 \(2\leqslant x<3\) 时，\([x]=2\)，所以
\[
C_x^8=\frac{8\cdot7}{x(x-1)}=\frac{56}{x(x-1)}
\]
由于 \(x(x-1)\) 在 \([2,3)\) 上递增，且在 \(x=2\) 时取值为 \(2\)，在 \(x\to3^-\) 时趋近于 \(6\)，故
\[
\frac{28}{3}<\frac{56}{x(x-1)}\leqslant28
\]
所以值域为
\[
\left(4,\frac{16}{3}\right]\cup\left(\frac{28}{3},28\right]
\]
故选 D.
\end{solution}

\section{填空题}

\begin{problem}
\(\lim\limits_{x \to 1} \frac{x - 1}{x^{2} + 3x - 4} =\)\fillinblank{}.
\end{problem}

\begin{answer}
\(\frac{1}{5}\)
\end{answer}

\begin{solution}
因式分解分母，得
\[
x^2+3x-4=(x-1)(x+4)
\]
于是
\[
\lim_{x\to1}\frac{x-1}{x^2+3x-4}
=\lim_{x\to1}\frac{1}{x+4}=\frac15
\]
\end{solution}

\begin{problem}
已知椭圆 \(\frac{x^2}{a^2} + \frac{y^2}{b^2} = 1\)（\(a > b > 0\)）的右焦点为 \(F\)，右准线为 \(l\)，离心率 \(e = \frac{\sqrt{5}}{5}\). 过顶点 \(A(0, b)\) 作 \(AM \perp l\)，垂足为 \(M\)，则直线 \(FM\) 的斜率等于\fillinblank{}.
\end{problem}

\begin{answer}
\(\frac{1}{2}\)
\end{answer}

\begin{solution}
设椭圆半焦距为 \(c\)，则 \(c^2=a^2-b^2\)，且
\(\frac ca=e=\frac{\sqrt5}{5}\)，\(b^2=a^2-c^2=\frac45a^2\)
从而 \(b=2c\). 右准线 \(l\) 的方程为 \(x=\frac{a^2}{c}\)，所以
\(M=\left(\frac{a^2}{c},b\right)\)，\(F=(c,0)\).
因此直线 \(FM\) 的斜率为
\[
k_{FM}=\frac{b}{\frac{a^2}{c}-c}
=\frac{b}{\frac{a^2-c^2}{c}}
=\frac{c}{b}=\frac12
\]
\end{solution}

\begin{problem}
设函数 \( y = f(x) \) 存在反函数 \( y = f^{-1}(x) \)，且函数 \( y = x - f(x) \) 的图象过点 \((1,2)\). 则函数 \( y = f^{-1}(x) - x \) 的图象一定过点\fillinblank{}.
\end{problem}

\begin{answer}
\((-1,2)\)
\end{answer}

\begin{solution}
由图象 \(y=x-f(x)\) 经过 \((1,2)\)，得
\(2=1-f(1)\)，\(f(1)=-1\).
所以函数 \(y=f(x)\) 的图象经过 \((1,-1)\). 反函数图象关于直线 \(y=x\) 对称，故 \(y=f^{-1}(x)\) 经过 \((-1,1)\). 在 \(x=-1\) 处，\(f^{-1}(x)-x=1-(-1)=2\)，所以所求点为 \((-1,2)\).
\end{solution}

\begin{problem}
已知函数 \(f(x) = \frac{\sqrt{3 - ax}}{a - 1}\)（\(a \neq 1\)）.
\begin{enumerate}
\item 若 \(a > 1\)，则 \(f(x)\) 的定义域是\fillinblank{}；
\item 若 \( f(x) \) 在区间 \( (0,1] \) 上是减函数，则实数 \(a\) 的取值范围是\fillinblank{}.
\end{enumerate}
\end{problem}

\begin{answer}
\(\left(-\infty,\frac{3}{a}\right]\)；\((-\infty,0)\cup(1,3]\).
\end{answer}

\begin{solution}
（1）当 \(a>1\) 时，分母为正，定义域由根式条件决定：
\[
3-ax\geqslant0\Longleftrightarrow x\leqslant\frac3a
\]
故定义域为 \(\left(-\infty,\frac3a\right]\).

（2）先讨论参数所在区间. 当 \(a>1\) 时，根式函数随 \(x\) 增大而减小，要使 \((0,1]\) 全在定义域内，必须有 \(3-a\geqslant0\)，即 \(1<a\leqslant3\)，此时除端点外严格减小.

当 \(0<a<1\) 时，根式随 \(x\) 增大而减小，而分母 \(a-1<0\)，所以 \(f(x)\) 在定义域内递增，不符合条件. 当 \(a=0\) 时，函数为常数，也不符合“减函数”的要求.

当 \(a<0\) 时，根式 \(\sqrt{3-ax}\) 随 \(x\) 严格增大，而分母为负，故 \(f(x)\) 在 \((0,1]\) 上严格减小. 因此
\[
a\in(-\infty,0)\cup(1,3]
\]
\end{solution}

\begin{problem}
对有 \(n\)（\(n \geqslant 4\)）个元素的总体 \(\{1, 2, 3, \cdots, n\}\) 进行抽样. 先将总体分成两个子总体 \(\{1, 2, \cdots, m\}\) 和 \(\{m + 1, m + 2, \cdots, n\}\)（\(m\) 是给定的正整数，且 \(2 \leqslant m \leqslant n - 2\)）. 再从每个子总体中各随机抽取 \(2\) 个元素组成样本，用 \(P_{ij}\) 表示元素 \(i\) 和 \(j\) 同时出现在样本中的概率，则 \(P_{1n} =\fillinblank{}\)；所有 \(P_{ij}\)（\(1 \leqslant i < j \leqslant n\)）的和等于\fillinblank{}.
\end{problem}

\begin{answer}
\(\frac{4}{m(n-m)}\)；\(6\).
\end{answer}

\begin{solution}
元素 \(1\) 在第一个子总体中，被抽中的概率为 \(\frac2m\)；元素 \(n\) 在第二个子总体中，被抽中的概率为 \(\frac2{n-m}\). 两个子总体的抽样相互独立，所以
\[
P_{1n}=\frac2m\cdot\frac2{n-m}=\frac{4}{m(n-m)}
\]

每次抽样从两个子总体各取 \(2\) 个元素，样本中恰有 \(4\) 个元素，因而样本内无序元素对的总数恒为
\[
\mathrm C_4^2=6
\]
若对每个样本中的元素对使用指示变量求和，再取数学期望，得到这些指示变量的期望之和正好是所有 \(P_{ij}\) 的和. 因此
\[
\sum_{1\leqslant i<j\leqslant n}P_{ij}=6
\]
\end{solution}

\section{解答题}

\begin{problem}
甲，乙，丙三人参加了一家公司的招聘面试，面试合格者可正式签约. 甲表示只要面试合格就签约. 乙，丙则约定：两人面试都合格就一同签约，否则两人都不签约. 设每人面试合格的概率都是 \(\frac{1}{2}\)，且面试是否合格互不影响. 求：
\begin{enumerate}
\item 至少有 \(1\) 人面试合格的概率；
\item 签约人数 \( \xi \) 的分布列和数学期望.
\end{enumerate}
\end{problem}

\begin{solution}
记甲、乙、丙面试合格分别为事件 \(\mathcal A\)、\(\mathcal B\)、\(\mathcal C\). 由题设，\(\mathcal A\)、\(\mathcal B\)、\(\mathcal C\) 相互独立，且
\[
P(\mathcal A)=P(\mathcal B)=P(\mathcal C)=\frac12
\]

（1）至少有 \(1\) 人面试合格的对立事件是三人都不合格，因此
\[
P(\text{至少有 }1\text{ 人合格})
=1-P(\overline{\mathcal A}\cap\overline{\mathcal B}\cap\overline{\mathcal C})
=1-\left(\frac12\right)^3
=\frac78
\]

（2）甲签约当且仅当 \(\mathcal A\) 发生；乙、丙签约当且仅当 \(\mathcal B\cap\mathcal C\) 发生，且二人同时签约. 因而
\[
\xi=\mathbf 1_{\mathcal A}+2\mathbf 1_{\mathcal B\cap\mathcal C}
\]
当 \(\mathcal A\) 不发生且 \(\mathcal B\cap\mathcal C\) 不发生时，\(\xi=0\)，所以
\[
P(\xi=0)=\frac12\left(1-\frac14\right)=\frac38
\]
当 \(\mathcal A\) 发生且 \(\mathcal B\cap\mathcal C\) 不发生时，\(\xi=1\)，所以
\[
P(\xi=1)=\frac12\left(1-\frac14\right)=\frac38
\]
当 \(\mathcal A\) 不发生且 \(\mathcal B\cap\mathcal C\) 发生时，\(\xi=2\)，所以
\[
P(\xi=2)=\frac12\cdot\frac14=\frac18
\]
当 \(\mathcal A\) 与 \(\mathcal B\cap\mathcal C\) 同时发生时，\(\xi=3\)，所以
\[
P(\xi=3)=\frac12\cdot\frac14=\frac18
\]
因此分布列为
\begin{center}
\begin{tabular}{c|cccc}
\(\xi\) & \(0\) & \(1\) & \(2\) & \(3\)\\ \hline
\(P\) & \(\frac38\) & \(\frac38\) & \(\frac18\) & \(\frac18\)
\end{tabular}
\end{center}
数学期望为
\[
E(\xi)=0\cdot\frac38+1\cdot\frac38+2\cdot\frac18+3\cdot\frac18=1
\]
\end{solution}

\begin{problem}
如图所示，四棱锥 \(P - ABCD\) 的底面 \(ABCD\) 是边长为 \(1\) 的菱形，\(\angle BCD = 60^{\circ}\)，\(E\) 是 \(CD\) 的中点，\(PA \perp\) 底面 \(ABCD\)，\(PA = 2\).
\begin{enumerate}
\item 证明：平面 \(PBE\) \(\perp\) 平面 \(PAB\)；
\item 求平面 \(PAD\) 和平面 \(PBE\) 所成二面角（锐角）的大小.
\end{enumerate}
\FigureLayoutDeclare{img/2008/hunan_science/q17_fig1.png}{7.0cm}{2bp 2bp 2bp 2bp}
\bitmapfigure[max width=0.36\linewidth]{img/2008/hunan_science/q17_fig1.png}
\end{problem}

\begin{solution}
建立直角坐标系，使底面在 \(z=0\) 内，并取
\(B=\left(0,0,0\right)\)，\(C=\left(1,0,0\right)\)，\(D=\left(\frac12,\frac{\sqrt3}{2},0\right)\)，\(A=\left(-\frac12,\frac{\sqrt3}{2},0\right)\).
由 \(PA\perp\) 底面且 \(PA=2\)，可取
\(P=\left(-\frac12,\frac{\sqrt3}{2},2\right)\)，\(E=\left(\frac34,\frac{\sqrt3}{4},0\right)\).

（1）有
\(\overrightarrow{BE}=\left(\frac34,\frac{\sqrt3}{4},0\right)\)，\(\overrightarrow{AB}=\left(\frac12,-\frac{\sqrt3}{2},0\right)\)
从而
\[
\overrightarrow{BE}\cdot\overrightarrow{AB}
=\frac38-\frac38=0
\]
又因为 \(BE\) 在底面内，而 \(PA\perp\) 底面，过 \(B\) 作直线 \(b\parallel PA\)，则 \(b\subset\) 平面 \(PAB\) 且 \(BE\perp b\). 直线 \(AB\) 与 \(b\) 在 \(B\) 相交，且 \(BE\perp AB\)，所以 \(BE\perp\) 平面 \(PAB\). 由于 \(BE\subset\) 平面 \(PBE\)，得到平面 \(PBE\perp\) 平面 \(PAB\).

（2）平面 \(PAD\) 的两个方向向量可取
\(\overrightarrow{AP}=\left(0,0,2\right)\)，\(\overrightarrow{AD}=\left(1,0,0\right)\)
故其一个法向量可取 \(\left(0,1,0\right)\). 平面 \(PBE\) 的一个法向量为
\[
\overrightarrow{BP}\times\overrightarrow{BE}
=\left(-\frac{\sqrt3}{2},\frac32,-\frac{\sqrt3}{2}\right)
\]
设两平面所成的锐角为 \(\theta\)，则
\[
\cos\theta
=\frac{\left\lvert\left(0,1,0\right)\cdot
\left(-\frac{\sqrt3}{2},\frac32,-\frac{\sqrt3}{2}\right)\right\rvert}
{\sqrt{\frac34+\frac94+\frac34}}
=\frac3{\sqrt{15}}
\]
由于 \(\theta\) 为锐角，
\[
\sin\theta=\sqrt{1-\frac9{15}}=\frac{\sqrt{10}}5
\]
所以所求二面角的大小为 \(\arcsin\frac{\sqrt{10}}5\).
\end{solution}

\begin{problem}
数列 \( \{a_{n}\} \) 满足 \(a_{1}=1\)，\(a_{2}=2\)，\(a_{n+2}=\left(1+\cos^{2}\frac{n\pi}{2}\right)a_{n}+\sin^{2}\frac{n\pi}{2}\)，\(n=1\)，\(2\)，\(3\)，\(\cdots.\)
\begin{enumerate}
\item 求 \(a_{3}\)，\(a_{4}\) 并求数列 \( \{a_{n}\} \) 的通项公式；
\item 设 \(b_{n} = \frac{a_{2n - 1}}{a_{2n}}\)，\(S_{n} = b_{1} + b_{2} + \dots + b_{n}\). 证明：当 \( n \geqslant 6 \) 时，\( |S_{n} - 2| < \frac{1}{n} \).
\end{enumerate}
\end{problem}

\begin{solution}
（1）当 \(n=1\) 时，\(\cos^2\frac{n\pi}{2}=0\)，\(\sin^2\frac{n\pi}{2}=1\)，所以
\[
a_3=a_1+1=2
\]
当 \(n=2\) 时，\(\cos^2\frac{n\pi}{2}=1\)，\(\sin^2\frac{n\pi}{2}=0\)，所以
\[
a_4=2a_2=4
\]

当 \(n=2k-1\) 时，递推式化为
\[
a_{2k+1}=a_{2k-1}+1
\]
奇数项构成首项为 \(a_1=1\)、公差为 \(1\) 的等差数列，故
\[
a_{2k-1}=k=\frac{2k-1+1}{2}
\]
当 \(n=2k\) 时，递推式化为
\[
a_{2k+2}=2a_{2k}
\]
偶数项构成首项为 \(a_2=2\)、公比为 \(2\) 的等比数列，故
\[
a_{2k}=2^k
\]
因此
\[
a_n=
\begin{cases}
\dfrac{n+1}{2},&n=2k-1,\ k\in\N^{*},\\
2^{n/2},&n=2k,\ k\in\N^{*}.
\end{cases}
\]

（2）由上式得
\(b_n=\frac{a_{2n-1}}{a_{2n}}=\frac{n}{2^n}\)，\(S_n=\sum_{k=1}^{n}\frac{k}{2^k}\).
将
\[
2S_n=1+\frac2{2}+\frac3{2^2}+\cdots+\frac n{2^{n-1}}
\]
与 \(S_n\) 相减，得到
\[
S_n=1+\sum_{k=1}^{n-1}\frac1{2^k}-\frac n{2^n}
=2-\frac{n+2}{2^n}
\]
所以
\[
\lvert S_n-2\rvert=\frac{n+2}{2^n}
\]
下面证明当 \(n\geqslant6\) 时 \(n(n+2)<2^n\). 当 \(n=6\) 时，\(6(6+2)=48<64=2^6\). 若对某个 \(r\geqslant6\) 有 \(r(r+2)<2^r\)，则
\[
2r(r+2)-(r+1)(r+3)=r^2-3>0
\]
从而
\[
(r+1)(r+3)<2r(r+2)<2^{r+1}
\]
由数学归纳法，\(n(n+2)<2^n\) 对一切 \(n\geqslant6\) 成立. 因此
\[
\lvert S_n-2\rvert=\frac{n+2}{2^n}<\frac1n
\]
\end{solution}

\begin{problem}
在一个特定时段内，以点 \(E\) 为中心的 \(7\text{ 海里}\) 以内海域被设为警戒水域. 点 \(E\) 正北 \(55\text{ 海里}\) 处有一个雷达观测站 \(A\). 某时刻测得一艘匀速直线行驶的船只位于点 \(A\) 北偏东 \(45^{\circ}\) 且与点 \(A\) 相距 \(40\sqrt{2}\text{ 海里}\) 的位置 \(B\)，经过 \(40\text{ 分钟}\) 又测得该船已行驶到点 \(A\) 北偏东 \(45^{\circ} + \theta\)（其中 \(\sin \theta = \frac{\sqrt{26}}{26}\)，\(0^{\circ} < \theta < 90^{\circ}\)）且与点 \(A\) 相距 \(10\sqrt{13}\text{ 海里}\) 的位置 \(C\). 求：
\begin{enumerate}
\item 求该船的行驶速度（单位：海里/小时）；
\item 若该船不改变航行方向继续行驶，判断它是否会进入警戒水域，并说明理由.
\end{enumerate}
\bitmapfigure[max width=0.20\linewidth]{img/2008/hunan_science/q19_fig1.png}
\end{problem}

\begin{solution}
建立平面直角坐标系，取点 \(A\) 为原点，正东方向为 \(x\) 轴正向，正北方向为 \(y\) 轴正向. 因为点 \(E\) 在点 \(A\) 正南方 \(55\) 海里处，所以
\[
E=(0,-55)
\]
点 \(B\) 在点 \(A\) 北偏东 \(45^{\circ}\) 且距离为 \(40\sqrt2\) 海里，故
\[
B=(40,40)
\]
由 \(0^{\circ}<\theta<90^{\circ}\) 及 \(\sin\theta=\frac{\sqrt{26}}{26}\)，得
\[
\cos\theta=\frac5{\sqrt{26}}
\]
于是
\[
\sin(45^{\circ}+\theta)
=\frac{\sqrt2}{2}\left(\cos\theta+\sin\theta\right)
=\frac3{\sqrt{13}}
\]
\[
\cos(45^{\circ}+\theta)
=\frac{\sqrt2}{2}\left(\cos\theta-\sin\theta\right)
=\frac2{\sqrt{13}}
\]
所以
\[
C=\left(10\sqrt{13}\sin(45^{\circ}+\theta),
10\sqrt{13}\cos(45^{\circ}+\theta)\right)=(30,20)
\]

（1）船在 \(40\) 分钟内行驶的路程为
\[
BC=\sqrt{(40-30)^2+(40-20)^2}=10\sqrt5\text{ 海里}
\]
由于 \(40\) 分钟为 \(\frac23\) 小时，船速为
\[
v=\frac{10\sqrt5}{2/3}=15\sqrt5\text{ 海里/小时}
\]

（2）直线 \(BC\) 的方程为
\[
2x-y-40=0
\]
其从 \(B\) 指向 \(C\) 的方向向量为 \(\overrightarrow{BC}=(-10,-20)\). 设点 \(E\) 到该直线的垂足为
\[
H=B+t\overrightarrow{BC}
\]
由 \(EH\perp BC\)，得
\[
t=\frac{(E-B)\cdot\overrightarrow{BC}}{\lvert\overrightarrow{BC}\rvert^2}
=\frac{(-40,-95)\cdot(-10,-20)}{500}
=\frac{23}{5}>1
\]
所以垂足 \(H\) 在船经过 \(C\) 后的航线上，且
\(H=\left(-6,-52\right)\)，\(EH=\sqrt{(-6)^2+3^2}=3\sqrt5<7\).
航线与以 \(E\) 为中心、半径为 \(7\) 海里的警戒圆的距离小于半径，故船会进入警戒水域.
\end{solution}

\begin{problem}
若 \(A\)，\(B\) 是抛物线 \(y^{2} = 4x\) 上的不同两点，弦 \(AB\)（不平行于 \(y\) 轴）的垂直平分线与 \(x\) 轴相交于点 \(P\)，则称弦 \(AB\) 是点 \(P\) 的一条“相关弦”. 已知当 \(x > 2\) 时，点 \(P(x, 0)\) 存在无穷多条“相关弦”. 给定 \(x_0 > 2\).
\begin{enumerate}
\item 证明：点 \(P(x_0, 0)\) 的所有“相关弦”的中点的横坐标相同；
\item 试问：点 \(P(x_0, 0)\) 的“相关弦”的弦长中是否存在最大值？若存在，求其最大值（用 \(x_0\) 表示）；若不存在，请说明理由.
\end{enumerate}
\end{problem}

\begin{solution}
（1）设点 \(P(x_0,0)\) 的任意一条相关弦的端点为
\(A=(x_1,y_1)\)，\(B=(x_2,y_2)\)，\(x_1\ne x_2\)
则 \(y_1^2=4x_1\)，\(y_2^2=4x_2\). 若 \(y_1+y_2=0\)，则 \(y_1^2=y_2^2\)，从而 \(x_1=x_2\)，与题设矛盾，所以 \(y_1+y_2\ne0\).

设弦中点为 \(M=(x_m,y_m)\)，则
\(x_m=\frac{x_1+x_2}{2}\)，\(y_m=\frac{y_1+y_2}{2}\).
由抛物线方程相减，得
\[
4(x_1-x_2)=y_1^2-y_2^2=(y_1-y_2)(y_1+y_2)
\]
因此弦 \(AB\) 的斜率为
\[
k=\frac{y_1-y_2}{x_1-x_2}=\frac4{y_1+y_2}=\frac2{y_m}
\]
其垂直平分线方程为
\[
y-y_m=-\frac{y_m}{2}(x-x_m)
\]
因为点 \(P(x_0,0)\) 在此直线上，所以
\[
-y_m=-\frac{y_m}{2}(x_0-x_m)
\]
又 \(y_m\ne0\)，故 \(x_m=x_0-2\). 这说明所有相关弦的中点横坐标都相同.

（2）由端点纵坐标的和与中点定义，结合 \(x_i=\frac{y_i^2}{4}\)，可得
\[
(y_1-y_2)^2=(y_1+y_2)^2-4y_1y_2
=4(4x_m-y_m^2)
\]
弦长 \(l\) 满足
\[
l^2=(x_1-x_2)^2+(y_1-y_2)^2
=\left(1+\frac{y_m^2}{4}\right)(y_1-y_2)^2
=(y_m^2+4)(4x_m-y_m^2)
\]
由端点不同及 \(y_m\ne0\)，有
\[
0<y_m^2<4x_m=4(x_0-2)
\]
令 \(t=y_m^2\)，则
\[
l^2=(t+4)\bigl(4x_0-8-t\bigr)
=4(x_0-1)^2-\bigl[t-2(x_0-3)\bigr]^2
\]
其中 \(0<t<4x_0-8\). 反过来，区间内每个 \(t\) 都给出两个不同的实数端点纵坐标，因而对应题设中的一条相关弦，所以这正是参数 \(t\) 的完整取值范围.

当 \(x_0>3\) 时，\(t_0=2(x_0-3)\) 满足
\[
0<t_0<4x_0-8
\]
故在 \(t=t_0\) 时取得最大值
\(l_{\max}^2=4(x_0-1)^2\)，\(l_{\max}=2(x_0-1)\).

当 \(2<x_0\leqslant3\) 时，\(2(x_0-3)\leqslant0\)，而 \(t>0\)，所以
\[
\frac{\mathrm d}{\mathrm dt}l^2=4(x_0-3)-2t<0
\]
此时 \(l^2\) 在允许的开区间上严格递减，最大的端点值只在 \(t=0\) 时才可能取得，但 \(t=0\) 对应 \(y_m=0\)，会使弦平行于 \(y\) 轴，已被题设排除. 因而不存在最大弦长.
\end{solution}

\begin{problem}
已知函数 \( f(x) = \ln^2 (1 + x) - \frac{x^2}{1 + x} \).
\begin{enumerate}
\item 求函数 \( f(x) \) 的单调区间；
\item 若不等式 \(\left(1 + \frac{1}{n}\right)^{n + \alpha} \leqslant \e\) 对任意的 \(n \in \mathbf{N}^*\) 都成立（其中 \(\e\) 是自然对数的底数），求 \(\alpha\) 的最大值.
\end{enumerate}
\end{problem}

\begin{solution}
函数定义域为 \((-1,+\infty)\)，且
\[
f'(x)=\frac{2\ln(1+x)}{1+x}-\frac{x^2+2x}{(1+x)^2}
=\frac{2(1+x)\ln(1+x)-x^2-2x}{(1+x)^2}
\]
令
\[
g(x)=2(1+x)\ln(1+x)-x^2-2x
\]
则
\[
g'(x)=2\bigl(\ln(1+x)-x\bigr)
\]
再令 \(h(x)=\ln(1+x)-x\)，则
\[
h'(x)=-\frac{x}{1+x}
\]
因此 \(h\) 在 \((-1,0)\) 上递增，在 \((0,+\infty)\) 上递减，且 \(h(0)=0\)，从而 \(h(x)<0\)（\(x\ne0\)）. 所以 \(g'(x)<0\)（\(x\ne0\)），\(g\) 在 \((-1,+\infty)\) 上递减. 又 \(g(0)=0\)，故
\(g(x)>0\)，\((-1<x<0)\)，\(g(x)<0\)，\((x>0)\).
由于分母 \((1+x)^2>0\)，得到 \(f'(x)>0\)（\(-1<x<0\)），\(f'(x)<0\)（\(x>0\)）. 因此 \(f(x)\) 在 \((-1,0)\) 上单调递增，在 \((0,+\infty)\) 上单调递减.

（2）原不等式等价于
\[
(n+\alpha)\ln\left(1+\frac1n\right)\leqslant1
\]
即
\(\alpha\leqslant G\left(\frac1n\right)\)，\(G(x)=\frac1{\ln(1+x)}-\frac1x\)，\(0<x\leqslant1\).
对 \(G\) 求导，得
\[
G'(x)=\frac1{x^2}-\frac1{(1+x)\ln^2(1+x)}
\]
由（1）知 \(f(x)<f(0)=0\)（\(x>0\)），即
\[
(1+x)\ln^2(1+x)<x^2
\]
所以 \(G'(x)<0\)，\(G\) 在 \((0,1]\) 上递减. 对任意正整数 \(n\)，有 \(\frac1n\leqslant1\)，于是
\[
G\left(\frac1n\right)\geqslant G(1)=\frac1{\ln2}-1
\]
因此 \(\alpha=\frac1{\ln2}-1\) 时所有 \(n\) 都满足原不等式，且 \(n=1\) 时给出 \(\alpha\leqslant G(1)\). 故 \(\alpha\) 的最大值为
\[
\frac1{\ln2}-1
\]
\end{solution}
